\documentclass[11pt]{article}
\usepackage[margin=1in]{geometry}
\usepackage[T1]{fontenc}
\usepackage[utf8]{inputenc}
\usepackage{amsmath,amssymb,amsthm}
\newtheorem{theorem}{Scoped theorem}
\setlength{\emergencystretch}{4em}
\setlength{\parindent}{0pt}
\setlength{\parskip}{0.6em}
\title{Bounded variation and the full bounded dual for analytic radial equations}
\author{Independent unverified proof draft}
\date{9 October 2026; uncompiled}
\begin{document}
\maketitle
Original P16.7 remains unresolved. All claims below retain the exact stated
analytic infrastructure and the explicit deeper published-source boundaries.
No incoming note supplies a correctness grant.
\begin{theorem}
For the bounded class $C_\theta$ of continuous locally uniformly ordinary
smooth middle-positive weights on a compact good smooth boundaryless analytic
surface, the genuine mixed products have conserved mass
$V=\int_X\theta^2$. With $\alpha_u=\theta+d'd''[u]$,
\[
 E(u)=\frac13\int_Xu(\theta^2+\theta\wedge\alpha_u+\alpha_u^2),
 \qquad \frac{d}{dt}E(u+t(v-u))=\int_X(v-u)\alpha_{u+t(v-u)}^2.
\]
The energy is concave and its tangent inequality holds for every bounded
competitor. Comparison holds for this entire bounded class.
\end{theorem}
\begin{theorem}
On the quotient over $K=\mathbf C((s))$, $|s|=e^{-1/8}$, with periods
$(s^8,2),(3,s^8)$, retain the specified published foundations. For
right-continuous nonincreasing $0\le h\le1$, constant $c$ on $[0,1)$,
put $Z=\int_0^\infty h^3<\infty$ and
$u_h=u_0-\int_0^{R}h$. Its genuine radial nonpluripolar measure has mass two:
\[
 \mu_h=(2-c^2)\delta_o+\xi_*[-d(h^2)],\qquad
 E(u_h)=-\frac16-\frac Z3.
\]
Its dual over the entire bounded class is exactly
\[
 \sup_{v\in C_\theta}\left(E(v)-\int_Xv\,d\mu_h\right)
 =\frac16+\frac{2Z}{3}.
\]
The unbounded operator and uniqueness asserted here are restricted to the
radial domain. No full intrinsic all-direction finite-energy theory follows.
\end{theorem}
Complete proofs, version distinctions, evidence and exact residual follow.

\section*{Full bounded dual on a concrete nonalgebraizable surface}

Unverified standalone scoped paper. Original P16.7: incomplete; candidate=false.

\section*{Bounded integration by parts, variation, and the tangent inequality}

This is an unverified proof draft. The theorem below concerns the entire specified bounded class, with all its ordinary analytic directions. It does not concern an unconstructed all-direction unbounded class.

Let X be a compact, good, separated, smooth, boundaryless analytic surface with its canonical CLD integration, and let theta be a smooth positive closed (1,1)-form with local smooth positive potentials q. Define C\_theta by requiring that u is continuous and that q+u, on a neighborhood of each point in the ordinary topology, is a locally uniform limit of ordinary smooth functions with middle-positive Hessian. Write DD=d'd'', alpha\_u=theta+DD[u], B(u,v)=alpha\_u wedge alpha\_v, and V=integral theta\textasciicircum{}2. Then the genuine local bounded products satisfy

E(u)=(1/3) integral u(theta\textasciicircum{}2+theta wedge alpha\_u+B(u,u)),
E(u+C)=E(u)+VC,
d/dt E(u+t(v-u))=integral(v-u)B(u+t(v-u),u+t(v-u)),
E(v)<=E(u)+integral(v-u)B(u,u).

Here all t in [0,1] are allowed, with the endpoint derivatives interpreted one-sided. In particular E is concave on all of C\_theta. We prove the integration by parts needed for these identities without assuming global smooth positive approximation.

The analytic infrastructure is exactly the ordinary current/form calculus of CLD, arXiv:1204.6277v2 (submitted 25 July 2025): compact smooth cutoffs and finite partitions in prescribed ordinary neighborhoods; density of smooth scalars on compact supports; canonical integration on a topologically separated pure two-dimensional boundaryless analytic space; symmetry and positive wedge calculus. The applicable statements are 6.5.2, 7.2.8/11/13, 7.7.1-8, 8.5.1/3/5, 8.7.1/2/4, and 15.7.1-2. Their hypotheses are retained. Only the smooth-factor Leibniz rule is used. The erroneous general continuous-multiplier closedness assertion 8.6.7 is not used.

First reconstruct the bounded products. In a moment presentation, symmetric (p,p) coefficients are symmetric matrices on the exterior coefficient space. Middle positivity is the PSD cone; the pth power of a Euclidean positive (1,1)-form is p! times its identity matrix. Thus every smooth symmetric test on a compact support is a difference of two compact middle-positive tests, using a cutoff and a finite partition. A positive symmetric current consequently has coefficient Radon measures: a scalar test times a positive form gives a positive functional bounded by the scalar sup norm, and Riesz representation applies. Independence of decompositions follows from equality on smooth tests and density. If positive T\_j converge weakly and b\_j converge locally uniformly, then b\_j T\_j converge to b T. Indeed choose a smooth g uniformly close to b and split the difference into (b\_j-g)T\_j+g(T\_j-T)+(g-b)T; the first and last terms are bounded by the positive test masses times the uniform errors.

Locally let p=q+u and r=q+v. Set B(u,v)=DD(p DD[r]). Simultaneous local smooth positive approximants, the preceding multiplication lemma, and smooth closed-factor Leibniz show that this is the limit of the positive smooth mixed products. It is symmetric, positive, closed, local, independent of approximants, and continuous when the potentials converge locally uniformly within the declared class. Smooth DD-null changes of q glue it. Globally B(u,v)=theta wedge alpha\_v+DD(u alpha\_v), and theta wedge alpha\_v=theta\textasciicircum{}2+DD(v theta). Testing the differentiated currents with the constant 1 gives integral B(u,v)=V. None of these assertions requires a global positive approximant. Finite convex combinations are in C\_theta by intersecting finitely many local approximation neighborhoods. The same argument and a smooth convex nondecreasing regularization prove finite-max stability.

Now fix a positive symmetric closed (1,1)-current T, held fixed throughout the next argument. Locally let f be a difference of two functions individually uniformly approximable by smooth positive weights. Choose smooth differences f\_i=p\_i-r\_i converging uniformly on compact subsets. The signed measure L\_T f=DD(fT) is the difference of the positive measures DD(pT) and DD(rT). For smooth p\_i, DD(p\_i T)=DD p\_i wedge T by smooth closed-factor Leibniz. Positivity, distributional convergence, and coefficient multiplication give L\_T f\_i -> L\_T f and f\_i L\_T f\_i -> f L\_T f, locally as currents; local signed variations are uniformly bounded. For example, if chi is a smooth nonnegative compact cutoff, then

integral chi DD p\_i wedge T=integral p\_i DD chi wedge T,

which is bounded uniformly because p\_i are locally uniformly bounded and DD chi wedge T is a fixed finite signed measure. A cutoff equal to one on a smaller compact set bounds the entire positive measure there.

For smooth a define Q\_T(a)=d'a wedge d''a wedge T. It is a positive measure; its bilinear polarization is symmetric because T annihilates antisymmetric (1,1)-tests. Smooth Leibniz gives

Q\_T(a)=(1/2)DD(a\textasciicircum{}2T)-a L\_T a,
integral chi Q\_T(a)=(1/2)integral a\textasciicircum{}2 DD chi wedge T-integral chi a L\_T a.

Apply this with a=f\_i-f\_j. If eta is its sup norm on a fixed compact neighborhood containing the support of chi, the positive decompositions of L\_T f\_i and L\_T f\_j just bounded give

0<=integral chi Q\_T(f\_i-f\_j)<=C\_chi(eta+eta\textasciicircum{}2).

The constant is independent of i,j. The same identity gives a uniform bound on integral chi Q\_T(f\_i). For any smooth compact (1,0)-test beta, coefficient PSD positivity gives Cauchy-Schwarz between beta and d''a, with chi=1 on the test support. This can be seen by applying T to (lambda beta+d'a) wedge J(lambda beta+d'a), integrating, and requiring the resulting quadratic to be nonnegative for every real lambda. The bound is the square root of the corresponding two positive quadratic masses. It is valid also for beta multiplied by any smooth bounded scalar.

Consequently d''f\_i wedge T are Cauchy when tested against beta, and, for another such difference g, f\_i d''g\_i wedge T are Cauchy: split their difference into (f\_i-f\_j)d''g\_i wedge T and f\_j d''(g\_i-g\_j) wedge T. The first is bounded by the uniform error in f times the bounded local gradient mass of g; the second tends to zero by the displayed estimate and the boundedness of f\_j. Interleaving any two approximation sequences proves independence. The resulting currents f d''g wedge T glue on overlaps, since the same estimates apply on every smaller compact subset of an overlap. This argument uses no positive approximation of T and no global approximation of f or g.

For smooth f\_i,g\_i, symmetry cancels the cross terms and yields

d'(f\_i d''g\_i wedge T-g\_i d''f\_i wedge T)
    =f\_i L\_T g\_i-g\_i L\_T f\_i.

Pass to the limit using the established convergence on each side. The local currents glue, so this is a global identity for global f,g of the stated local difference class. On compact X, test with 1 to obtain the required symmetry

integral f DD(gT)=integral g DD(fT).

Also the local limits of Q\_T(f\_i) equal Q\_T(f)=(1/2)DD(f\textasciicircum{}2T)-f DD(fT), a positive top measure. In particular

integral f DD(fT)=-integral Q\_T(f)<=0.

Notice that f\textasciicircum{}2T is merely a well-defined coefficient product. It is never asserted to be closed, nor is f\textasciicircum{}2 asserted to be psh. This is why the proof does not invoke 8.6.7.

Every u in C\_theta is locally the difference (q+u)-q; hence these identities apply to global u,v and to T=theta or alpha\_w for any w in C\_theta. They also apply by subtraction to their signed differences. The identification DD((u-v)alpha\_w)=B(u,w)-B(v,w) follows directly from the reconstructed recursive products.

Put f=v-u, u\_t=u+tf and alpha\_t=alpha\_\{u\_t\}. Difference multilinearity makes E(u\_t) a polynomial of degree at most three. Differentiate its three terms and integrate by parts with T=theta and alpha\_t:

E'(u\_t)=(1/3) integral f[theta\textasciicircum{}2+theta alpha\_t+alpha\_t\textasciicircum{}2+(alpha\_t-theta)(theta+2alpha\_t)]
        =integral f alpha\_t\textasciicircum{}2.

There are no limiting derivative assertions in this polynomial calculation. Differentiating again gives

E''(u\_t)=2 integral f DD(f alpha\_t)=-2 integral Q\_\{alpha\_t\}(f)<=0.

Thus the derivative decreases along the segment. Integrating it from 0 to 1 proves the tangent inequality for every u,v in the entire C\_theta. Constants cancel in its dual objective when a measure has raw mass V. This closes the bounded first-variation/concavity gate under the stated analytic current infrastructure, without solving an unbounded all-direction gate.

There is also comparison for the ENTIRE bounded class. For epsilon>0 set w=max(u,v-epsilon), which belongs to C\_theta by the local convex regularization already proved. Locality on its two strict open regions and positivity imply

V=integral MA(w)>=integral\_\{u>v-epsilon\} MA(u)+integral\_\{u<v-epsilon\} MA(v).

Subtracting the first term from V yields integral\_\{u<v-epsilon\} MA(v)<=integral\_\{u<=v-epsilon\} MA(u). As epsilon decreases to zero, both sets increase to \{u<v\}. Continuity from below for finite measures proves integral\_\{u<v\} MA(v)<=integral\_\{u<v\} MA(u). Possible mass on each intermediate equality set causes no difficulty, because it was included on the right and its increasing union is exactly the final strict set. No global approximation or unbounded comparison theorem is used.

Artifacts: n1.md, P16.7-bounded-dual-paper.tex, proof-companion.md and source-application-ledger.md. Rerun the own integrity and finite algebra diagnostics with `python identity\_audit.py` and `python diagnostics.py`. Neither program establishes the universal analytic argument above. Support: none.

\section*{Concrete quotient, calibrated branch products, and foundation boundaries}

Retain exactly K=C((s)), |s|=exp(-1/8), t=s\textasciicircum{}8, and X=(G\_m\textasciicircum{}2)\textasciicircum{}an/<(t,2),(3,t)>. The following reconstructs the concrete inputs for applying n1 and the radial formulas. All claims are unverified drafts under the explicitly named published analytic foundations, never grants from supplied notes.

In coordinates x\_j=-log|z\_j| the deck group translates x by Z\textasciicircum{}2. Ordinary inverse images of rectangles of widths less than one have disjoint nonidentity translates. They give analytic quotient charts. The relation is a locally finite union of closed translation graphs: compact tropical coordinate ranges bound every possible deck translation. Conrad-Temkin, Non-archimedean analytification of algebraic spaces, author version 24 February 2009, Lemma 4.2.3 glues quotients of restrictions to an ordinary open cover, with open images covering the quotient. Lemma 4.2.4 gives coherent G-sheaf descent along a flat quasi-finite surjection. Its immediate proof reduces to finite faithfully flat affinoid algebras, where completed tensor products become ordinary tensor products. Descending the coherent ideal of the relation proves that the quotient diagonal is closed. The source's entire applied statements and immediate arguments were read in the retained primary bodies; deeper Berkovich descent results remain external foundations.

For that descent choose a FINITE subcover of the injective quotient charts on the compact tropical fundamental square, and let Z be their disjoint union. Each chart maps by an open immersion, so Z->X is flat, surjective and quasi-finite, with each fibre bounded by the number of charts; the same holds for Z times Z->X times X. The pulled-back diagonal is locally a finite disjoint union of closed translation graphs and its coherent ideal has the evident graph cocycle. Lemma 4.2.4 descends this ideal. We do not incorrectly apply quasi-finite descent to the infinite-sheeted torus->X map.

The strict compact polyannulus x\textasciicircum{}\{-1\}([0,1]\textasciicircum{}2) surjects onto X. Centers in (1/8 Z)\textasciicircum{}2 within distance 1/16 of any tropical lift, and rectangles of width 1/2 inside width 3/4 quotient charts, give strict affinoid neighborhoods. Local torus charts give smoothness, pure dimension two, and empty relative boundary. Compactness, separatedness and empty boundary give properness. The nested integral Laurent polyannuli in the torus form a connected exhaustion. Thus X is good, strict, smooth, separated, proper, boundaryless and connected of dimension two over this base K.

For completeness the algebraic exclusion has an explicit foundation boundary. Gruson, Fibres vectoriels sur un polydisque ultrametrique (1968), IV.5 Theorem 1 gives SK\textasciicircum{}0(H)=0 for finite completed tensor products of his disc algebras T\_\{1,r\}, r in Q tensor v(K*), and unit-circle algebras. No wide annulus is an input to that theorem. A rank-one projective P has [P]=[H], so P direct-sum Q is isomorphic to H direct-sum Q for some projective Q; determinants cancel and trivialize P. Radii in |K*| rescale to unit factors. Induct on the number of wide annuli. For A(r,R) times Y, trivialize a line bundle on C(r) times Y, glue it to the trivial bundle on D(r) times Y, and obtain a line bundle on D(R) times Y. Both the boundary and filled space have one fewer wide factor. The domains D(r) and A(r,R) are a finite rational G-cover of D(R), intersecting in C(r): use |z/a|<=1 and |a/z|<=1, with |a|=r and (a,z) the unit ideal. Tate-Kiehl coherent descent, in Van der Put, Rigid analytic spaces, GAU J7 (1975-76), 2.14-2.16.3, justifies this gluing. Its complete printed J7-09 descent statement/diagram was inspected; its deeper Kiehl proof is a published foundation. Hence every finite product in question, including P\_m=\{|s|\textasciicircum{}m<=|z\_j|<=|s|\textasciicircum{}\{-m\}\}, has Picard group zero.

For a unit sum b\_beta z\textasciicircum{}beta on P\_m, the Gauss log norm and that of its inverse are convex and sum to zero, so are affine. A maximizing monomial at an interior point must give that affine function everywhere. Its exponent alpha gives |b\_beta/b\_alpha| exp((m/8)|beta-alpha|\_1)<=1. Thus the unit is b\_alpha z\textasciicircum{}alpha(1+g), with zero constant coefficient and ||g||\_\{P\_j\}<=|s|\textasciicircum{}\{m-j\} for j<m. A global unit on the whole torus is a scalar monomial, since another exponent would dominate on a ray. Normalize consecutive frames by scalar monomials so e\_\{m+1\}|P\_m=(1+g\_m)e\_m. The products p\_j=product\_\{m=j\}\textasciicircum{}infinity(1+g\_m)|P\_j converge to units, including inverse tails. Frames e\_j p\_j are compatible and glue on the interior exhaustion. Hence Pic((G\_m\textasciicircum{}2)\textasciicircum{}an)=0.

Commuting line-bundle multipliers f(lambda\_i z)=c\_i z\textasciicircum{}\{h\_i\}f(z) forces, for h\_1=(A,B), h\_2=(C,D), t\textasciicircum{}C2\textasciicircum{}D=3\textasciicircum{}At\textasciicircum{}B. Valuations give C=B and prime factorization gives A=D=0. Thus h\_1=(0,k), h\_2=(k,0). If k!=0 take n=(1,sign k), lambda=lambda(n), h=(kn\_2,kn\_1), so Q=v(lambda\textasciicircum{}h)=2|k|>0. For a nonzero Laurent coefficient a\_m, the exact recurrence is a\_\{m+jh\}=a\_m c\textasciicircum{}j lambda\textasciicircum{}\{-jm-hj(j+1)/2\}. Its valuation has negative quadratic leading term -Qj(j+1)/2, contradicting Laurent convergence on the unit polyannulus. If k=0, pairing the difference of two exponents with both period valuations shows they coincide. Every nonzero section is a single nowhere-zero monomial.

For any local meromorphic fraction F=g/h, its intrinsic pole ideal J=\{b:bF is analytic\} is the coherent kernel of multiplication by g into O/(h). Regular affinoid localizations are UFDs (Stacks 0AG0); cancellation makes J locally free of rank one. Flat affinoid restriction and coherent descent glue these ideals. The normal height-one intersection theorem (Stacks 031T) also verifies pole cancellation. A global frame of J on the torus gives analytic f\_0 and f\_1=Ff\_0 with common deck multipliers. The preceding recurrence makes their ratio constant. Thus Gamma(X,M)=K; no general analytic fraction equalizer is asserted.

A separated finite-type algebraic-space algebraization Y would be smooth of dimension two by Conrad-Temkin 2.3.1 and 2.3.2. The former reflects separatedness, monomorphism, surjectivity, isomorphism, open immersion, flatness, smoothness and etaleness, and for finite-type morphisms also closed immersion, finiteness, properness and quasi-finiteness. Connectedness follows by analytifying clopen components; regularity makes the irreducible components disjoint, so Y is integral. A dense scheme locus (Stacks 086U) supplies its generic field of transcendence degree two. Stacks 0ENA/0END/0EN6 apply: regular noetherian etale charts have no embedded weak associates, satisfy the finite-component conditions 0BB1, and separatedness gives decency and field monomorphisms at generic points. They identify global meromorphic functions with K(Y). Coherent analytification and the regular algebraic-to-analytic local maps of Ducros, Families of Berkovich spaces, 2.1.4 inject K(Y) into Gamma(X,M). That item requires a scheme locally of finite type over an affinoid algebra; our scheme charts have base K. It cites Ducros 2009 Theorem 3.3 and Berkovich 1993 Proposition 2.6.2 for its deeper regularity/flatness. Flat local maps preserve nonzerodivisors, so numerator/denominator sections analytify; a closed point avoiding their zeros detects every nonzero rational function. The injection into K contradicts transcendence degree two. This excludes all separated finite-type K-algebraic spaces at base K, with no arbitrary-extension assertion. The exact whole-source verification requirement of the assignment is not certified completed; in particular all book contexts and deeper descent/flatness trees have not been read. This argument's applicability is therefore stated under these published foundations rather than as an unconditional foundational audit.

Now put q=(x\_1\textasciicircum{}2+x\_2\textasciicircum{}2)/2 and theta=DD q. Deck changes of q are affine, so theta descends. The canonical identity-torus calibration has saturated lattice Z\textasciicircum{}2 and degree one, giving theta\textasciicircum{}2=2 i\_*dx for probability Haar dx. Raw mass is two. With G(a)=a\textasciicircum{}2/2-a/2+1/12 periodically extended and u0=-G(x\_1)-G(x\_2), convolution gives A=alpha\_\{u0\}, A\textasciicircum{}2=2delta\_o and theta wedge A=i\_*(da tensor delta\_0+delta\_0 tensor da), o=i(0). Since G has mean zero, its value at zero is 1/12 and its minimum is -1/24, the fibre F=tau\textasciicircum{}\{-1\}(0) has u0=-1/6, E(u0)=-1/6 and sup u0=1/12.

On the ordinary sheet |x\_j|<1/4 put w\_j=z\_j-1 and W=\{|w\_j|<1\}. Define R=-log max(|w\_1|,|w\_2|) on W, zero off W, with its sole infinite value at y=(1,1). This is continuous into [0,infinity]. For every real N>0 let D\_N=-min(N,R). If ell\_j=log max(exp(-N),|w\_j|) and b\_j=log max(1,|z\_j|), then D\_N=max(ell\_1-b\_1,ell\_2-b\_2) on the sheet, and

q+u0+D\_N=max(ell\_1+b\_2,b\_1+ell\_2)-(1/2)log|z\_1z\_2|-1/6.

The maximum is convex in ordinary unit moments. Near a zero omit its inactive capped logarithm. Softmax gives the required local positive smooth approximants, so U\_N=u0+D\_N belongs to precisely C\_theta. Both theta and A vanish on W, and D\_N vanishes off W; coefficient multiplication gives D\_N theta=D\_N A=0.

For a finite nonincreasing step profile h=sum a\_j 1\_\{[0,N\_j)\}, 0<N\_1<...<N\_k, a\_j>=0 and c=sum a\_j<=1, set g(t)=integral\_0\textasciicircum{}t h and u\_h=u0-g(R). Convexity gives membership in C\_theta. On W its weight, up to a constant, is f(y\_1,y\_2)=sum a\_j max(-N\_j,y\_1,y\_2). For c>0 its softmax Hessian is a positive covariance and its gradient is a diffeomorphism onto the interior of the triangle p\_1,p\_2>0, p\_1+p\_2<c. Injectivity follows from strict convexity. For each interior p the softmax minus p dot y is coercive, since c max(0,y\_1,y\_2)-p dot y is positive in every nonzero direction; minimization gives surjectivity and the positive Hessian gives a smooth inverse. The uniform error is at most epsilon c log 3. For h(N\_j)<p\_1+p\_2<h(N\_j-), the limiting minimizer is uniquely (-N\_j,-N\_j): positive coordinates of p force both y coordinates to their maximum, and the remaining scalar slopes straddle p\_1+p\_2. Uniform error and coercivity imply convergence of inverse gradient points. Exceptional slope boundaries have area zero.

Change variables in 2 det Hess f\_epsilon. The area of the corresponding triangle annular band is (h(N\_j-)\textasciicircum{}2-h(N\_j)\textasciicircum{}2)/2. Dominated convergence therefore gives mass h(N\_j-)\textasciicircum{}2-h(N\_j)\textasciicircum{}2 at the actual translated Gauss point xi\_\{N\_j\}. The analytic Gauss section is continuous; composing arbitrary continuous analytic tests with it proves this measure identity on the analytic carrier, with the exterior factor 2 and actual degree one retained. Away from vertices the local weight is affine or depends on one moment after an affine subtraction; at zero coordinates omit the inactive logarithm, so there is no other branch mass.

The outside mass is at o. On F with exactly one |w\_j|<1 the relevant weights depend on at most one ordinary unit moment and have zero square. Where |z\_1|=|w\_1|=|z\_2|=|w\_2|=1, use the common moment (z\_1,w\_1,z\_2,w\_2). CLD-v2 8.8.6 supports tropical difference-of-convex top products on that moment's top skeleton. CLD 12.2.7-10 bounds local tropical dimension by the graded residual transcendence degree of these coordinates. Its essential input is now read directly: Ducros, Espaces de Berkovich, polytopes, squelettes et theorie des modeles, Confluentes Math. 4(4) (2012), Theorem 3.4. For a Gamma-strict k-analytic space and invertible analytic f\_1,...,f\_n, it supplies a compact Gamma-strict analytic neighborhood U of x whose tropical germ is unchanged on compact analytic sub-neighborhoods; its dimension is at most the transcendence degree over graded k of the field generated by the reduced coordinates, with pure equality when the germ is boundaryless. Our strict quotient chart satisfies this with Gamma=\{1\}. Its entire immediate proof uses Theorem 1.3 and Temkin germ reduction; these deeper dependencies are explicitly retained as published inputs. If the reductions of z\_1,z\_2 are dependent, the degree is at most one, so there is no top mass near the point. If independent, reducing a Laurent polynomial after dividing by a maximal-norm coefficient shows its norm is exactly the original unit Gauss norm. This identifies the point as o. Total mass two now proves

MA(u\_h)=(2-c\textasciicircum{}2)delta\_o+xi\_*[-d(h\textasciicircum{}2)].

Zero profiles are handled directly by A\textasciicircum{}2. The three energy integrals are 0, -1/6 and -1/3-integral h\textasciicircum{}3: the first uses mean zero; the second uses D theta=D A=0; the third uses u0=-1/6 on F, mass two, and Tonelli integral g[-d(h\textasciicircum{}2)]=integral h\textasciicircum{}3. Hence E(u\_h)=-1/6-(1/3)integral h\textasciicircum{}3. In particular MA(U\_N)=delta\_o+delta\_xi\_N for all real N>0.

Artifacts and reruns: n2.md and the standalone paper contain this reconstruction; `python identity\_audit.py` retains the applied source bodies/pixels under sources-LOCALONLY, and `python diagnostics.py` checks only finite rational formula instances. No supplied result is a premise; the explicitly named published analytic foundations remain essential. Support: n1.

\section*{Continuous Stieltjes radial equations and the full bounded dual}

Use the actual quotient, bounded calculus and calibrated finite formulas reconstructed in n1-n2. The new result is an upper bound, and exact value, for the dual over EVERY v in the entire bounded C\_theta. The constructed unbounded operator remains radial. No all-direction E1 recognition or full uniqueness is asserted.

Let H consist of right-continuous nonincreasing profiles h with 0<=h<=1, constant c on [0,1), and Z(h)=integral\_0\textasciicircum{}infinity h\textasciicircum{}3<infinity. Put g\_h(r)=integral\_0\textasciicircum{}r h, u\_h=u0-g\_h(R), with its canonical value at R=infinity. Then h tends to zero. Let nu\_h=-d(h\textasciicircum{}2) be the Stieltjes measure on [1,infinity), including a jump at 1. Its mass is c\textasciicircum{}2 and nu\_h((t,infinity))=h(t)\textasciicircum{}2 for t>=1. The statements proved below are

MA(u\_h)=(2-c\textasciicircum{}2)delta\_o+xi\_*nu\_h=:mu\_h,
B(u\_h,u\_k)=(2-cd)delta\_o+xi\_*[-d(hk)],
E(u\_h+C)=-1/6-Z(h)/3+2C,
sup\_\{v in C\_theta\}(E(v)-integral v mu\_h)=1/6+2Z(h)/3.

The first two measures are genuine locally bounded analytic products off y extended by zero at y, and also genuine full-mass nonpluripolar truncation limits. The last supremum is the FULL bounded dual, rather than a pencil or radial dual.

Choose compact finite step h\_m from below: on [0,1) keep c, on nested dyadic partitions of [1,2\textasciicircum{}m) sample h at right endpoints, and put zero afterward. Refining the grid and enlarging the interval makes h\_m increase to h at continuity points (indeed with right-continuous representatives everywhere at finite positive depths). Their primitives g\_m increase to g\_h, uniformly on each finite interval because integral\_0\textasciicircum{}R |h-h\_m| tends to zero. Thus u\_m=u0-g\_m(R) decrease to u\_h, locally uniformly on X minus \{y\}. The boundedness of R on a compact subset of that complement follows because the closed bidiscs \{R>=r\} decrease to \{y\}; compactness gives a finite depth bound.

On W minus \{y\}, put a=max(log|w\_1|,log|w\_2|)<0. The local weight is phi(a)-1/6, where phi(a)=-g\_h(-a). Its derivative is h(-a), nonnegative and nondecreasing in a, so phi is convex and nondecreasing. Composing with the maximum is convex. Near a zero of one coordinate the other has nonzero norm and the inactive coordinate is omitted. Convex smooth approximation on ordinary unit moment neighborhoods proves membership in the declared local bounded class. Outside W, the ordinary open \{R<1\} has u\_h=u0+cD\_1. Consequently the genuine local operator is already available everywhere off y; it is not assigned by an inverse measure convention.

By n2 the actual bounded measures are mu\_m=(2-c\textasciicircum{}2)delta\_o+xi\_*[-d(h\_m\textasciicircum{}2)]. Stieltjes tails converge on bounded depth intervals, with the usual tests at continuity endpoints. At a continuity depth r their tail above r is at most h(r)\textasciicircum{}2, tending to zero as r tends to infinity, uniformly in m. Finite partition approximation of a continuous test on [1,r], followed by discarding this tail, proves weak convergence on the depth line. Composing the test with the continuous actual analytic section xi gives weak convergence against every continuous analytic test on X. The finite initial mass and outside part converge exactly. The resulting pushforward is Radon: nu\_h is finite and tight, compact depth intervals have compact continuous images in the Hausdorff X, and its tail can be discarded. The branch mass is c\textasciicircum{}2; there is no atom or escaped mass at y. Local product continuity off y identifies this weak measure with the already constructed analytic local product. This proves the first formula and total mass two. Polarize the finite formulas and pass locally uniformly and weakly for h\_m k\_m to get the second formula. Its branch masses are nonnegative since hk decreases, their total is cd, and the outside mass is 2-cd. Both tails vanish. These arguments give independence of the described approximants and any others converging locally uniformly with initial values and uniformly tight tails converging as stated.

The full-mass truncation recipe for continuous profiles requires an endpoint check. Take M>0. If g\_h(infinity)<=M, then u\_h is bounded; integral |h-h\_m| tends to zero on the whole depth line, so u\_m converge uniformly on X and its measure agrees with the bounded operator. If M<g\_h(infinity), continuity and monotonicity give a unique finite positive r\_M with g\_h(r\_M)=M. A plateau at level M cannot occur in this case: a nonnegative nonincreasing derivative that reaches zero stays zero, so such a plateau would have value g\_h(infinity). Set h\textasciicircum{}M(t)=h(t) for t<r\_M and zero for t>=r\_M. Then g\_\{h\textasciicircum{}M\}=min(g\_h,M) and

u\textasciicircum{}M=u0-min(g\_h(R),M)=max(u\_h,u0-M)

is in C\_theta. If r\_M<1 this profile is outside H but is an allowed bounded profile: n2 proved all positive real cutoffs. Approximate h\textasciicircum{}M by finite step profiles with positive real cutoffs and integrate; the convergence is uniform globally since the support lies in [0,r\_M]. On each ordinary neighborhood the same convex moment presentation, or a finite local max regularization, supplies smooth positive approximants. Near y the potential is exactly u0-M. This is a direct membership proof, independent of a global approximation theorem.

Its branch Stieltjes measure is nu\_h restricted to depths below r\_M, with the entire endpoint mass replaced by h(r\_M-)\textasciicircum{}2 delta\_\{r\_M\}. This includes any original jump at that endpoint. For r\_M<1 the endpoint mass is c\textasciicircum{}2. The outside part stays (2-c\textasciicircum{}2)delta\_o. The strict open set \{g\_h(R)<M\} is \{R<r\_M\}; it excludes the entire clamp atom, including the original endpoint atom. Bounded locality therefore gives

1\_\{g\_h(R)<M\} MA(u\textasciicircum{}M)=1\_\{g\_h(R)<M\} mu\_h.

As M increases these restricted measures increase to mu\_h off y; their omitted tail is h(r\_M-)\textasciicircum{}2, which tends to zero when g\_h(infinity)=infinity. Thus their mass tends to two. In the bounded case sufficiently large M already gives the bounded operator. This proves the genuine nonpluripolar truncation recipe. For sufficiently large integers j, max(u\_h+C,-j) differs from u\_h+C only on W, where u0=-1/6, and equals the same clamp with M=j+C-1/6. Hence ordinary constant truncations give the identical recipe. For mixed products, clamp each slot and restrict to the intersection of the two strict unclamped sets. Locality gives the local B there; the omitted branch tail is bounded by h(r)k(r) at a common large depth and tends to zero. This proves mixed full mass as well, with no rigid atom.

The first energy integral is integral u\_h theta\textasciicircum{}2=0. The second is integral u\_h(theta wedge alpha\_\{u\_h\})=-1/6: the difference from u0 annihilates theta and A locally, exactly as in n2, and these reference measures are supported away from y. For the top measure u0=-1/6 on its support F and Tonelli gives

integral g\_k(t) nu\_h(dt)=integral\_0\textasciicircum{}infinity k(t)h(t)\textasciicircum{}2 dt,
integral g\_h(t) nu\_h(dt)=Z(h).

These are integrals of nonnegative functions, so there is no infinity cancellation. The top energy integral is -1/3-Z(h). Thus the displayed energy is the genuine three-term energy in this domain, and also the limit of the bounded energies. Constants add twice their value.

For all h,k in H, Holder gives integral k h\textasciicircum{}2<=||k||\_3 ||h||\_3\textasciicircum{}2. It therefore justifies signed subtraction and differentiation of the cubic along every radial segment. With h\_t=(1-t)h+tk and constants C\_t=(1-t)C+tD,

d/dt E(u\_\{h\_t\}+C\_t)=-integral(k-h)h\_t\textasciicircum{}2+2(D-C)
                         =integral[(u\_k+D)-(u\_h+C)] mu\_\{h\_t\},
d\textasciicircum{}2/dt\textasciicircum{}2 E(u\_\{h\_t\}+C\_t)=-2 integral(k-h)\textasciicircum{}2 h\_t<=0.

The possible undefined difference of two infinite values at y is assigned zero on that null singleton. This proves every segment first variation and concavity WITHIN the radial class. It also proves energy monotonicity under pointwise order there, by integrating the first variation against a positive measure.

Decreasing continuity with variable constants has its own proof. Suppose u\_\{h\_m\}+C\_m decreases pointwise and its energies are bounded below. Reading the potential at a fixed point outside F shows C\_m decreases. The bound E\_m<=-1/6+2C\_m bounds C\_m below; hence C\_m decreases to finite C. Put delta\_m=C\_m-C. Then g\_m-delta\_m increases pointwise, is uniformly 1-Lipschitz and concave, and is bounded on each finite interval. Its limit is a finite concave primitive g with g(0)=0, constant derivative on [0,1), and derivative h\_m -> h almost everywhere by concave secant inequalities. The energy bound and bounded constants bound Z\_m; Fatou gives Z(h)<infinity. Primitive convergence is locally uniform. The canonical value at infinity is correct because sup\_m sup\_r(g\_m(r)-delta\_m)=sup\_r sup\_m(g\_m(r)-delta\_m).

To prove energy convergence, order gives g\_m<=g+delta\_m. For a nonnegative nonincreasing compact test a, Stieltjes integration by parts gives integral a h\_m<=integral a h+delta\_m a(0). Take a=h\_m\textasciicircum{}2 1\_\{[0,L)\} and let L tend to infinity. Since a(0)<=1, Holder gives

Z\_m<=Z(h)\textasciicircum{}\{1/3\} Z\_m\textasciicircum{}\{2/3\}+delta\_m.

Boundedness and this inequality yield limsup Z\_m<=Z(h); Fatou supplies the opposite liminf. Thus energies converge, including variable constants. Monotonicity gives r h\_m(r)\textasciicircum{}3<=Z\_m, so the common cubic bound gives uniformly vanishing squared measure tails. Stieltjes convergence and the preceding tightness argument yield weak MA continuity; two such slots yield mixed continuity by the product tail bound. No general all-direction decreasing theorem is inferred.

Comparison with variable constants is also intrinsic to this radial domain. Let f=g\_h-g\_k and b=C-D. Inside W the strict comparison set is the open depth set \{f>b\}; outside W it is present exactly if b<0. On a component interval (a,l) with positive finite contact endpoints, the right derivative at a is h(a)-k(a)>=0 and the left derivative at l is h(l-)-k(l-)<=0. Therefore the difference of branch masses on that OPEN interval equals

[h(a)\textasciicircum{}2-k(a)\textasciicircum{}2]-[h(l-)\textasciicircum{}2-k(l-)\textasciicircum{}2]>=0.

An infinite right endpoint contributes zero tails. For an initial interval at depth zero, if b=0 the same initial derivative sign holds. If b<0 its initial term c\textasciicircum{}2-d\textasciicircum{}2 need not be positive, but the included outside mass difference is d\textasciicircum{}2-c\textasciicircum{}2, exactly cancelling it. If b>0 there is no initial interval. Summing the countably many intervals proves mu\_k\{u\_h+C<u\_k+D\}<=mu\_h\{u\_h+C<u\_k+D\}. Open endpoints exclude their atoms, which explains the right/left tail convention. Both measures give y mass zero. This is not comparison with arbitrary potentials outside the radial domain.

For fixed L, \{h in H:Z(h)<=L\}, with C=0, is compact in the metric sum\_\{m>=1\}2\textasciicircum{}\{-m\}min(1,sup\_\{[0,m]\}|g\_h-g\_k|). Indeed the primitives are concave, nondecreasing, 1-Lipschitz and zero at zero. Diagonal extraction on rational depths gives uniform convergence on finite intervals. Secant inequalities give almost-everywhere derivative convergence, and Fatou preserves Z<=L. Choose the right-continuous representative of the limiting derivative. This proves compactness in a metrizable topology. It is separated: finite Gauss depth segments detect g, the initial linear segment is detected by g(1), and the canonical value at y is determined by g(infinity). This is precisely locally uniform convergence on compact subsets off y: depth is bounded on such compacts, and conversely the compact analytic Gauss segments detect uniform primitive convergence. All normalized u\_h have sup u\_h=1/12, attained outside F. The topology does not make the value at y continuous; it fixes that value canonically.

Every positive Radon nu on [1,infinity) of mass c\textasciicircum{}2<=1 with integral\_1\textasciicircum{}infinity nu((t,infinity))\textasciicircum{}\{3/2\} dt finite has a genuine radial solution: take h=c on [0,1), h(t)=sqrt(nu((t,infinity))) afterward. The Stieltjes measure -d(h\textasciicircum{}2) is exactly nu, including its atom at 1, so the genuine constructed operator is (2-c\textasciicircum{}2)delta\_o+xi\_*nu. Conversely its tails recover h uniquely, so uniqueness modulo constants holds within this radial class. Discrete depth grids are the subcase of step tails. The literal finite-energy condition is Z=c\textasciicircum{}3+integral\_1\textasciicircum{}infinity nu((t,infinity))\textasciicircum{}\{3/2\}; the branch mass bound belongs to this image family alone.

Now prove the new FULL bounded dual statement. For the finite-step u\_m already constructed, apply n1's tangent inequality to every v in C\_theta:

E(v)-integral v mu\_m<=E(u\_m)-integral u\_m mu\_m=1/6+(2/3)Z(h\_m).

Here integral u\_m mu\_m=-1/3-Z(h\_m), with the same raw mass two throughout. Weak convergence mu\_m -> mu\_h against the continuous v and Z(h\_m) increasing to Z(h) yield

E(v)-integral v mu\_h<=1/6+(2/3)Z(h).

This is an upper bound for the ENTIRE C\_theta, not only radial competitors. For the reverse inequality test u\_m against the ACTUAL limiting analytic measure. Tonelli gives integral u\_m mu\_h=-1/3-integral h\_m h\textasciicircum{}2, so

E(u\_m)-integral u\_m mu\_h=1/6-Z(h\_m)/3+integral h\_m h\textasciicircum{}2
                         ->1/6+2Z(h)/3.

Monotone convergence applies since h\_m<=h. The two inequalities give exact equality and full bounded-dual admissibility of EVERY finite-energy continuous or discrete radial solution measure in this range. If Z is infinite for a profile decreasing to zero, the same lower tests diverge; one should not regard that measure as finite-energy eligible. For bounded profiles, u\_h itself attains the full bounded dual. For unbounded profiles the proof only needs the bounded optimizing sequence.

For example h=1 on [0,N), zero afterward gives mu=delta\_o+delta\_xi\_N and full bounded dual 1/6+2N/3. This is distinct from the measure 2delta\_xi\_N, for which 1/6+5N/3 is still only a restricted lower bound. The present theorem neither proves its full admissibility nor supplies a universal atom cap. Nor does bounded-dual finiteness here recognize membership in a full intrinsic all-direction E1 class: no such class and extension are proved.

A genuinely unbounded example is h=1/i on [i-1,i), i=1,2,... . Then g(infinity)=sum\_i 1/i=infinity, Z=sum\_i 1/i\textasciicircum{}3<infinity, and the actual measure is delta\_o+sum\_\{i>=1\}(1/i\textasciicircum{}2-1/(i+1)\textasciicircum{}2)delta\_\{xi\_i\}. The positive branch weights telescope to one. The strict nonpluripolar clamp construction above gives full raw mass two and no atom at the rigid y; its full bounded dual is 1/6+(2/3)sum\_i 1/i\textasciicircum{}3. This provides an explicit unbounded solution and its full bounded admissibility, with uniqueness only in the stated radial domain.

Artifacts: n3.md, P16.7-bounded-dual-paper.tex and proof-companion.md contain the entire argument; `python diagnostics.py` checks finite tangent-gap and cubic identities, and `python build\_deliverables.py` reproduces document/schema evidence. Mathematical correctness is the universal argument, not execution. Support: n1, n2.

\section*{Independent audit results, restricted duals, and the exact original residual}

This is an unverified scoped paper, not a resolution of original P16.7. Notes n1-n3 supply reconstructed proofs rather than incoming correctness grants. In particular the new full bounded tangent inequality is what upgrades the constructed radial measures from restricted-dual finiteness to full C\_theta-dual finiteness. The operator on the unbounded class in n3 remains a genuinely constructed radial operator.

We first independently check the ROOT384 distinction between its two restricted test families. Let mu be a positive measure of raw mass two supported at actual branch points xi\_N, N>=1; let T(t)=mu\{xi\_N:N>t\}. Thus T is nonincreasing, 0<=T<=2, and T=2 on [0,1). For a finite nonincreasing profile 0<=h<=1, supported in a bounded interval and constant on [0,1), n2 and Tonelli give

E(u\_h)-integral u\_h mu=1/6+integral\_0\textasciicircum{}infinity [h(t)T(t)-h(t)\textasciicircum{}3/3]dt.

All measures and test functions here are actual analytic objects. This calculation only evaluates a restricted supremum; it says nothing about an upper bound on the full C\_theta.

For each fixed T>=0 the maximum over 0<=a<=1 of aT-a\textasciicircum{}3/3 is Phi(T), with maximizing a=min(1,sqrt(T)), where

Phi(T)=(2/3)T\textasciicircum{}(3/2) for 0<=T<=1,
Phi(T)=T-1/3 for 1<=T<=2.

The optimizer is itself nonincreasing as a function of depth. Truncate it at large depths and approximate it from below by nested finite steps as in n3. Its objective integrand increases: on a<=min(1,sqrt(T)), the derivative T-a\textasciicircum{}2 is nonnegative. Monotone convergence proves that the exact finite-combination radial restricted dual is 1/6+integral Phi(T), finite exactly when the tail T belongs to L\textasciicircum{}(3/2). This is an exhaustive analytic optimization argument, not a finite experiment.

The smaller pencil consists of u0+aD\_N, 0<=a<=1, N>=1. Put A(N)=integral\_0\textasciicircum{}N T. Its objective is 1/6+aA(N)-a\textasciicircum{}3N/3. Maximizing a gives 1/6+(2/3)A(N)\textasciicircum{}(3/2)/N\textasciicircum{}(1/2) if A(N)<=N, and 1/6+A(N)-N/3 if A(N)>=N. Since T tends to zero for a finite depth measure, A(N)/N tends to zero. Hence this pencil dual is finite exactly when sup\_\{N>=1\} A(N)/N\textasciicircum{}(1/3)<infinity. The elementary inequalities N T(N)<=A(N) and integral\_1\textasciicircum{}N t\textasciicircum{}(-2/3)dt=3(N\textasciicircum{}(1/3)-1) show that this is equivalent to T(N)<=C N\textasciicircum{}(-2/3). This is the weak-tail criterion and is strictly weaker than L\textasciicircum{}(3/2).

For an actual critical countable example take N\_i=2\textasciicircum{}i, i=0,1,..., and branch masses 2(1-2\textasciicircum{}(-beta))2\textasciicircum{}(-beta i), with beta=2/3. They telescope to two. Their depth tail is bounded above and below by positive multiples of t\textasciicircum{}(-2/3) on t>=1. The pencil criterion holds, whereas each dyadic interval contributes a fixed positive amount to integral T\textasciicircum{}(3/2), so the finite-combination radial dual is infinite. The full bounded dual is at least that supremum and is therefore infinite. This measure is excluded by energy admissibility and cannot be an eligible counterexample to original P16.7. For beta>2/3 these restricted duals are finite; this argument does not prove full C\_theta-dual finiteness for that total-branch-mass-two family.

For mu=2delta\_\{xi\_N0\}, T=2 on [0,N0) and zero afterward. The optimizer h=1 on that interval gives exact restricted value 1/6+5N0/3. This is still only a LOWER bound on the full bounded dual. By contrast n3's actually solved measure delta\_o+delta\_\{xi\_N0\} has the exact FULL bounded dual 1/6+2N0/3. The outside atom and branch weights matter. Failure of the chosen radial family to realize branch mass two is not a universal mass-cap theorem and is not a counterexample to the literal construction axis.

We also check the positivity convention used in the bounded proof. On a real four-dimensional moment space, the coefficient space for (2,2)-forms is Sym\textasciicircum{}2(Lambda\textasciicircum{}2 R\textasciicircum{}4), of dimension 21. Strong generators are ww\textasciicircum{}T with w decomposable. In Plucker coordinates write w=(x,y) in R\textasciicircum{}3 direct-sum R\textasciicircum{}3, with decomposability x dot y=0. A quadratic form vanishing on every such w has pure x and pure y parts zero, by taking y=0 or x=0. Write its cross part x\textasciicircum{}T B y. Its vanishing whenever x is perpendicular to y implies B\textasciicircum{}T x is parallel to x for every x. Applying this to basis vectors and their pairwise sums forces B to be scalar. Thus the annihilator of the strong span is exactly one-dimensional, and the strong span is 20-dimensional. One must not claim it is the entire 21-dimensional coefficient space. The middle PSD cone DOES span the entire symmetric coefficient space and has the Euclidean identity as an order unit. That is the cone used in n1. None of this says that multiplying a closed current by an arbitrary continuous function preserves closedness.

The independent whole-mathematics audit covered the ROOT384 proof body and its binding/scope body, ROOT385 proof body and binding/scope body, the retained whole383 proof companion, and native272 whole387 companion. It separately read all five actual native272 FINAL note texts; those are not the five expanded note texts. The expanded bodies were read in their exact entirety inside the companion, with separately read summaries and support fields. The supplied ledger and preamble were read, and the complete Markdown-to-TeX ordered encoding was checked mechanically. Hash identity was used to identify duplicated physical content, never counted as an additional fresh mathematical reading. The exact coverage and incomplete source/book coverage are recorded in source-application-ledger.md. No supplied program was executed or imported.

The main scoped verdicts are these. The finite branch formulas, raw mass two, all-three-integral energy, continuous Stieltjes extension, clamp endpoints, variable-constant decreasing energy argument, radial comparison and separated radial compactness have reconstructed proofs in n2-n3. The source-dependent outside atom is justified by the actual local-dimension theorem read in Ducros 2012 Theorem 3.4, not just metadata for CLD reference [50]. Its immediate proof is read and its deeper Temkin/valuation dependencies remain named published foundations. Full proof credit for Fresnel-Van der Put 1983, KST 2016, or deeper Conrad-Temkin/Ducros foundations is not asserted. Gerritzen 1977's actual full OCR and printed 3.1-3.2 pixels give the H\textasciicircum{}1 corrections, not arbitrary affinoid Picard invariance. Gruson's applied theorem is only for pure discs and unit circles; wide annuli are handled by the separately proved finite rational hole-filling induction. The alleged gaps are repaired or explicitly bounded, not replaced by an undifferentiated Open verdict.

For reproducible observations, the own standard-library integrity program checked all 757 bound records (183 mathematics, 557 primary evidence, 17 book context), with no byte or SHA256 mismatch. The five actual-versus-expanded native texts have five distinct SHA256 pairs; every expanded text is an exact substring of the whole companion. Every nonempty Markdown line occurs in order, with character escaping, in the 67551-byte TeX encoding. These are identity facts, not correctness facts. The own finite algebra diagnostic enumerated 35 rational monotone three-step profiles and all 1225 ordered pairs. It verified the radial tangent gap

(1/3)sum length\_i (k\_i-h\_i)\textasciicircum{}2(k\_i+2h\_i)>=0

and the cumulative-majorization inequality in those finite instances. For h=(5/6,1/2,0) at depths (1,3,7), it found branch masses (4/9,1/4,0), outside mass 47/36, Z=179/216, energy -287/648, potential integral -251/216 and dual value 233/324; all are exact fractions. The clamp observations include M=1/4 at depth 3/10 with clamp atom 25/36, and M=4/3 at depth 2 with clamp atom 1/4. These observations only check the algebra illustrated by the universal proofs. The runtime was Python 3.12.14 on Linux 6.18.44 x86\_64 with glibc 2.41. No resource-heavy computation was launched; no supervised facility was supplied or needed for these bounded diagnostics. The requested native model selection is not independently inspectable from this role, and no additional native/model call was made.

Original P16.7 still requires a finite-energy operator and corresponding existence and uniqueness theorem on genuinely nonalgebraizable analytic spaces, with an explicit positivity datum and replacements for compactness, envelopes, energy differentiation and comparison. This work supplies the datum theta on a concrete nonalgebraizable surface under the named published foundations, a full bounded theory and a rigorous unbounded radial equation domain. It does NOT construct an intrinsic all-direction E1 space with approximation-independent mixed products, nonpluripolar full mass and decreasing continuity for every finite-energy potential. It does not establish separated compactness of that full space, obstacle envelopes and orthogonality for all analytic obstacles, full unbounded first variation, a sufficient criterion for EVERY arbitrary-support measure, surjectivity for every such admissible measure, or full uniqueness modulo constants. Radial compactness and radial uniqueness cannot replace those requirements. The fixed 2delta\_\{xi\_N0\} full-admissibility and universal solution/obstruction questions remain unresolved. All seventeen whole book chapters and every deeper primary proof tree have not been independently read in their entirety, so the mandatory unrestricted whole-source audit is not certified complete either.

An attempted shortcut would define all-direction E1 solely as decreasing limits of bounded C\_theta objects and invoke n1. Its missing step is precise: n1 proves continuity only for locally uniform bounded convergence. A decreasing bounded sequence can fail to converge locally uniformly near its singular set; uniform coefficient multiplication estimates then do not provide mixed-product convergence, preservation of mass, or approximation independence. The radial clamp proof controls this singularity explicitly by the vanishing Stieltjes tail. No corresponding all-direction control has been obtained. Finite bounded dual is therefore not used as an unproved recognition theorem for all-direction E1. The source book's enlargement of test classes likewise uses decreasing regularization and energy continuity within the original BFJ signature, which are exactly the missing transfers here.

Scoped verdict: the bounded integration-by-parts/variation/concavity/tangent/comparison gates and full bounded-dual admissibility for the genuine finite-energy radial solutions are proved as drafts under the stated analytic foundations. Original verdict: incomplete; candidate=false. No acceptance, count increase, native compilation, publication, new tab or protected-editor mutation has occurred.

Artifacts: n1.md through n4.md, P16.7-bounded-dual-paper.tex, proof-companion.md, source-application-ledger.md, final-notes.json, scoped-verdict.json, original-verdict.json, evidence/packet-identities.json, evidence/version-identities.json and evidence/diagnostics.json. Rerun only the OWN programs with `python identity\_audit.py`, `python diagnostics.py`, and `python build\_deliverables.py`. All source copies are LOCALONLY and their filenames are listed in the ledger. Support: n1, n2, n3.

\section*{Fresh reading and application ledger}

This ledger distinguishes mathematical reading, application checking, machine identity checking, and retained material. All incoming files, ledgers, source records, notes, and programs were DATA. No incoming note supplied correctness credit. Only the own notes n1, n2, and n3 are used as local supports. Source-dependent steps use the named primary statements below, with their literal hypotheses checked. All source copies and book contexts remain LOCALONLY. No supplied/native program was executed or imported, no extra agent or model call was made, and no compiler, acceptance, publication or protected-editor action was performed.

The requirement to read every deeper primary proof and every whole book context has NOT been fulfilled. This is an explicit unfinished audit obligation. The results below are scoped proofs under named published foundations, rather than a claim that the unrestricted foundational audit or original P16.7 has been completed.

\subsection*{Entire mathematical proof bodies and version distinctions}

Packet indices refer to the exact three supplied binding arrays, copied into inputs-LOCALONLY. Every input identity was checked before an emission. The reader logs an emission, not an automatic reading event; accidental truncated emissions were not credited and were repaired with smaller/overlapping reads. Physical duplicate content was identified and not counted again as fresh reading.

Mathematics 137: whole ROOT384 proof, 16162 bytes. Final untruncated character intervals [0,8500) and [8500,end) covered the entire body after a combined emission had been truncated. Reconstructed applications: actual finite convex weights, strict covariance and gradient polygon, degree-one Gauss section, all three integrals, pencil weak-tail optimization, finite-combination L\textasciicircum{}(3/2) optimization, critical excluded-energy example. The exact finite-combination and pencil optimizations are re-proved in n4; n2-n3 re-prove analytic product and energy ingredients.

Mathematics 138: entire ROOT384 scope/binding body, 4303 bytes. Read as data; no historical correctness grant used. Its lower-bound status for 2delta\_\{xi\_N\} is maintained.

Mathematics 139: whole ROOT385 proof, 14879 bytes, untruncated. Mathematics 140: entire ROOT385 binding/scope body, 3715 bytes. The continuous Stieltjes formula, mixed products, every radial segment, variable-constant decreasing majorization and energy convergence, endpoint comparison and separated off-rigid compactness are reconstructed in n3. They do not establish all-direction E1.

Mathematics 134: entire retained whole383 companion, 66573 characters, including its complete proof and 21055-byte ledger. Fresh intervals [0,23000), repaired by [19000,36000), then [36000,52000) and [52000,end) after the initial output truncation. The reconstructed proof in n2 retains base K, normalization, all separated finite-type algebraic-space exclusion, disc/circle Gruson restriction, rational annulus filling, compatible frames, all-signed Laurent recurrence and intrinsic pole ideal. It explicitly identifies the deeper source boundary.

Mathematics 136: entire qualified\_whole\_reading field, 7696 characters, including ROOT383 bindings and scope. Large arrays were machine identity checked, not given fresh proof-reading credit.

Mathematics 171: entire native272 companion, 63788 characters, emitted in three complete chunks [0,21000), [21000,42000), [42000,end). These contain all five EXPANDED text bodies and the entire 11827-byte source ledger. Mathematics 181 is that exact ledger substring, not a second fresh reading.

Mathematics 146: entire actual native272 FINAL JSON, 20306 bytes. All five ACTUAL text bodies were read. Mathematics 148: five EXPANDED notes JSON, 55975 bytes. The summaries and supports were separately read; each complete expanded text is an exact substring already physically read in 171. All five expanded text bodies differ from their actual FINAL counterparts. Their individual actual/expanded byte lengths are respectively (5060,25004), (2732,6515), (1315,3165), (5069,11274), and (2184,5951). These sizes refer to text fields; the JSON totals include summaries, supports, and escaping. This work never substitutes expanded text for what the actual FINAL claimed.

Mathematics 142: complete Markdown is the exact suffix of the already read companion 171. Mathematics 143: 67551-byte TeX. Every nonempty Markdown line occurs, correctly escaped and in order. All remaining nonempty TeX segments were separately inspected: there are exactly two, the preamble/title/introductory original-false sentence and the final end-document. Thus the mathematical body is an identified complete encoding, not a hash-only claim or a separate fresh duplicate reading. See evidence/TeX-unique-segments.json.

Mathematics 156: entire exact historical ROOT387 scope/binding report, 4906 bytes. Mathematics 182: entire qualified\_whole\_reading field, 4905 characters; the arrays are integrity evidence. Mathematics 141: entire ROOT Gerritzen correction record, 4188 bytes. Mathematics 179: entire scoped verdict, 967 bytes. The entire README/checks/diagnostic/environment/version/ROOT-identity reports at 144,173,174,175,176,177,180 were read as reports, never mathematical correctness grants.

The independently checked corrections are substantive. The outside atom in the actual radial measure is delta\_o; finite profiles have outside mass 2-c\textasciicircum{}2. Actual real-depth U\_N have MA=delta\_o+delta\_\{xi\_N\}. Continuous Stieltjes and discrete countable scopes agree by the construction in n3. The new tangent bound is over all C\_theta, whereas prior 2delta\_\{xi\_N\} restricted values remain lower bounds. Native whole387 does not assert a universal branch cap; neither does this work.

\subsection*{Applied primary statements and actual reading}

CLD: Chambert-Loir and Ducros, arXiv:1204.6277v2, 25 July 2025. Full extract and PDF retained at primary 4 and 5. Fresh selected character ranges of the authentic extract include [74100,79000), [446153,453500), [481000,485200), [523794,531000), [610050,616500), [668089,687100), and [1446200,1450500). Complete applied literal source units 327 (8.8.6, 6885 bytes), 328 (12.2.7-10, 4823 bytes), and 329 (5.2.5, 7704 bytes) were read. They give positive order-zero/Radon calculus, ordinary smooth cutoff/form integration, smooth closed-factor Leibniz, recursive local continuity, calibrated analytic integration and support, and tropical local dimension. The final adjacent partial statement in the 1.2 range was not given full-proof credit. The full 2MB extract was retained but not wholly read; deeper canonical-integration proof trees are published background, not newly verified full trees.

Applicability: all n1 spaces are good, topologically separated, pure dimension two and without boundary, and products use local ordinary smooth-positive approximability. Moment multiplicities are degree one with saturated tangent lattice Z\textasciicircum{}2; the top exterior factor is two. Coefficient control uses the full middle PSD cone, not the 20-dimensional strong span. Ordinary smooth positivity is retained, not replaced by weak harmonic smoothness or a sheaf-generated stage-recognition claim. Only smooth closed-factor Leibniz is used. The source's continuous-multiplier closedness claim 8.6.7 is rejected as a premise; n1 gives its own Dirichlet-current construction instead. No global positive approximation theorem is invoked.

Ducros: Espaces de Berkovich, polytopes, squelettes et theorie des modeles, Confluentes Math. 4(4), 2012, DOI 10.1142/S1793744212500077. Authentic full extract/PDF at primary 6/7. Metadata/title checked in [0,3500); ENTIRE Theorem 3.4 and its immediate proof were read in [129000,143000), printed pages 41-43. Literal hypotheses: Gamma-strict k-analytic space, invertible local coordinates f\_i; the stabilized tropical germ has dimension bounded by graded residual transcendence degree, with pure equality at a boundaryless germ. Our strict torus quotient charts permit Gamma=\{1\}. Where z\_j and w\_j=z\_j-1 are units the generated residual field is at most one-dimensional unless z\_1,z\_2 reduce independently; independence forces the unit Gauss norm. This closes the applied [50] local-dimension gate, rather than using bibliographic metadata. Its Theorem 1.3 and Temkin germ-reduction foundations were NOT wholly re-proved or wholly audited; the conditional foundation boundary is explicit. Full 57-page PDF retained, not claimed wholly freshly read.

Gruson 1968: full extract/PDF retained at primary 22/23; actual entire applied IV.5 theorem and immediate proof unit at primary 120, 10037 bytes, fresh reading. Its finite completed tensor products are disc algebras T\_\{1,r\}, r in the rationalized value group, and pure unit circles. SK\textasciicircum{}0=0 and rank-one determinant cancellation suffice. No arbitrary wide annulus is claimed in that source. The own finite rational G-cover annulus-hole filling induction in n2 makes the required extension.

Van der Put, Rigid analytic spaces, GAU J7 (1975-76): entire applied coherent-descent unit 121, 6242 bytes, fresh reading. Printed J7-09 pixel 320 was actually viewed to repair omitted OCR formulas and inspect the descent diagram. The finite-cover gluing includes the cocycle compatibility. Deeper Tate/Kiehl descent is a named published foundation, not a full newly audited proof tree.

Gerritzen, Zerlegungen der Picard-Gruppe nichtarchimedischer holomorpher Raume, Compositio Math. 35 (1977), 23-38: complete ROOT OCR extract primary 316 was read in [18000,end) and [0,18000), covering the full 30516-character body; full PDF retained at 315. Printed section 3.1-3.2 pixels 321,322,323 were actually viewed; formulas and positive/negative Laurent unit factorization were checked. The correct decompositions are Pic(X times E)=Pic(X) direct-sum H\textasciicircum{}1(X,G\_1) and Pic(X times L)=Pic(X) direct-sum H\textasciicircum{}1(X,G\_1) direct-sum H\textasciicircum{}1(X,G\_\{-1\}) direct-sum H\textasciicircum{}1(X,Z\_X). They are not arbitrary affinoid Picard invariance. The Laurent elimination proof has nth largest positive coefficient removed by multiplication with 1-f\_0\textasciicircum{}\{-1\}f\_n xi\textasciicircum{}n and norm error bounded by the square of the previous error; successive eliminations and the convergent unit product give the factorization. Deeper older units/classification references are named background. This source is used as an audit correction, not needed instead of n2's Picard proof.

Fresnel-Van der Put 1983 at primary 313/314 and KST 2016 August version at 317/318: whole files retained, no fresh full-proof credit and no theorem from their whole bodies used as an unproved substitute. In particular pro-Picard, contractibility and a formal-model claim do not replace the actual annulus induction.

Conrad-Temkin, Non-archimedean analytification of algebraic spaces, author 24 February 2009: complete applied body primary 243, 23302 characters, read in [0,11000), [11000,end); primary 244, 24257 characters, read in [0,12500), [12500,end). These include 2.3.1/2.3.2 reflection of smoothness and other properties and 4.2.3 quotient gluing/4.2.4 coherent flat quasi-finite descent with full immediate units. Coherent extension body 247 was read through character 12500, including its applied coherent analytification/site construction; the remainder was not claimed whole. Lemma 4.2.4 is applied to a FINITE injective chart cover Z->X and its product, not the infinite-sheeted torus covering. The chart map is flat, surjective and quasi-finite with finite bounded fibres. Deeper Berkovich descent facts remain published background. Direct raw PDF retrieval failed in the supplied native history; no new retrieval attempt or extra native call occurred here. That failure does not supply proof credit; the retained complete applied primary units are the actual source evidence used.

Ducros, Families of Berkovich spaces, 2.1.3/2.1.4: complete immediate units in primary 243 were read. The algebraic-to-analytic local map is regular, hence flat, for schemes locally of finite type over an affinoid algebra; our algebraic scheme charts have base K. Its cited Ducros 2009 Theorem 3.3 and Berkovich 1993 Proposition 2.6.2 remain external published foundations, not newly fully read/re-proved. Flatness preserves nonzerodivisors and permits the meromorphic injection in n2. No claim about arbitrary analytic fraction equalizers is used.

Stacks applied statements: complete primary 248, 15839 characters, read for 0AG0, 0FJH and 031T including immediate proofs. The exact meromorphic-functions hypotheses and proofs in primary 24 units 6-10 (0ENA, 0END, 0EN6, 0BB1, 086U) were read in full; 0EN6 also occurs in the complete read primary 243. Regular noetherian etale charts have no embedded weak associates, finite irreducible-component conditions, and separatedness supplies decency and generic field monomorphisms. The dense scheme locus for a decent space and the integral meromorphic field identify K(Y). These are exact applied published statements, not a claim to have newly audited every recursive Stacks dependency.

\subsection*{Book contexts and original residual}

All seventeen whole book contexts were bound-byte checked and retained under book-context-LOCALONLY. They were NOT all freshly read. Fresh book passages, read only as mathematical context, are chapter 16 lines 2331-2605 (differences and local energy currents), chapter 28 lines 594-795 (test class enlargement and dual at MA), and chapter 72 lines 5500-5610 (the original family and residual boundaries). Inventories of chapters 20,16,28,72 were examined; inventories are not whole mathematical readings. No chapter result is taken as a support premise in n1-n3. Chapter 28 makes the exact missing transfer plain: enlarging bounded tests to every intrinsic finite-energy potential requires decreasing regularization and energy continuity in the relevant theory, which are not established by the bounded tangent inequality alone.

The original P16.7 statement remains exactly the nonalgebraizable analytic construction/equation axis, not a narrowed radial target. Retained gates: K=C((s)), |s|=exp(-1/8), periods (s\textasciicircum{}8,2),(3,s\textasciicircum{}8), good/strict/smooth/proper/boundaryless/connected dimension two and exclusion of all separated finite-type K-algebraic spaces. The current paper supplies explicit theta, bounded mixed products and full bounded variation/tangent/comparison, genuine radial nonpluripolar finite-energy equations, and exact full bounded dual admissibility for their measures. It does not prove all-direction intrinsic E1/full mass/decreasing continuity/separated compactness, all analytic obstacle envelopes and orthogonality, all-direction unbounded variation, every arbitrary-support admissibility criterion and equation, or full uniqueness. The fixed 2delta\_\{xi\_N\} question is unresolved. No universal local atom cap or original counterexample is claimed.

\subsection*{Reproduction and evidence classes}

Run `python identity\_audit.py`, `python diagnostics.py`, `python build\_deliverables.py` from the working directory. The reader can be rerun with `python audit\_io.py math 137 --start 0 --end 8500`, or any documented packet/index/interval, with JSON selections `--unit`, `--block`, `--key`; book line selections use `--lines FIRST LAST`. Every program is own standard-library code. Sources and supplied programs are never executed/imported. None of these commands compiles TeX or communicates externally.

evidence/packet-identities.json: all 757 identity rows match. evidence/version-identities.json: all five actual/expanded distinctions, ledger and Markdown identities and complete ordered TeX encoding. evidence/TeX-unique-segments.json: the only two nonencoded segments, read explicitly. emissions.jsonl: actual emitted intervals; it explicitly warns against interpreting emission as reading/correctness. evidence/diagnostics.json: exact finite rational observations, 35 profiles and 1225 pairs, not proof. evidence/environment.json: Python/platform/configuration limits. evidence/build-checks.json: own serialization and body-preservation checks, not mathematical acceptance.

The standalone deliverable is P16.7-bounded-dual-paper.tex with every complete own note and this ledger encoded; proof-companion.md is the readable counterpart. final-notes.json and actual-final-notes.json contain the same own full note texts. There is no separate expanded version of our result. scoped-verdict.json and original-verdict.json preserve conditional source boundaries, the concrete new gates, and candidate=false. ROOT remains responsible for any compiler, independent correctness review, acceptance or publication.

\end{document}
